% !TEX program = xelatex
% this file is encoded in utf-8
% v1.7

\documentclass[12pt, a4paper]{ntust_report} 

\usepackage{amssymb}
%\usepackage{libertine}
%\usepackage[T1]{fontenc}
%\usepackage{lmodern}
%\usepackage{kpfonts}
%\usepackage[sc]{mathpazo}
%\usepackage{microtype}
%\usepackage{slantsc}% http://ctan.org/pkg/slantsc

\usepackage{fontspec}	% To enable font settings 加這個就可以設定字體 
\usepackage{xeCJK}      % To let Chinese-English font set in separate location 讓中英文字體分開設置

%\usepackage[indentfirst=true]{xeCJK}
%\setmainfont{Arial}            %設定主要字型，也就是英文字型 
\setmainfont[Mapping=tex-text]{Times New Roman}
%\setmainfont{Linux Libertine O}
%\setmainfont{Junicode}
%\setCJKmainfont{DFKai-SB}     %設定中文字型 
%\setCJKmainfont{BiauKai}

\setCJKmainfont{AR PL KaitiM Big5}
\XeTeXlinebreaklocale "zh"      		% These two lines have to be added, 這兩行一定要加，中文才能自動換行 
\XeTeXlinebreakskip = 0pt plus 1pt      % 	so that the Chinese sentence can change line automatically.



% Unless school modifies the thesis format (margins, header, footer, watermark)
% or needs to improve the use of LaTeX packages,
% or needs to change the default fonts or the encoding,
% then no need to modify the content inside this file.
% For thesis writing personalization, please modify the files starting with 'my_'

\usepackage[nospace]{cite}  		% for smart citation
\usepackage{geometry}  				% for easy margin settings
%\usepackage{subfigure}  			% for subfigure
%\usepackage[dvipdfm]{graphicx}  	% for graphic using eps
\usepackage[xetex]{graphicx}
%\usepackage{graphicx}  			% for graphic using eps

%\usepackage{epstopdf} 				% for automatic transformation (eps to pdf) while using pdflatex


%\usepackage{algorithmic}  			% for algorithm
\usepackage{algorithm}
\usepackage{algorithmicx}
\usepackage{algpseudocode}

%
% margins setting
\geometry{verbose,a4paper,tmargin=3.5cm,bmargin=2cm,lmargin=3cm,rmargin=3cm}
%
\usepackage{amsmath} 				% for mathematical function (AMS Format)
%\usepackage{amssymb} 				% for mathematical symbol (AMS Format)
										% if 'amssymb' and 'Xunicode' conflict, please put 'amssymb' beforehand
\usepackage{mathrsfs} 				% for cursive mathematical symbol,
										% for example, typing '\mathscr{E} may result 'cursive E'
\usepackage{listings} 				% for listing any programming code
%
% listing setting
\lstset{breaklines=true,			% a long line of code can be broken into lines of code
extendedchars=false,				
texcl=true,							
comment=[l]\%\%,					
basicstyle=\small,					% approximately 10 pt
commentstyle=\upshape,				% the default is 'italics', so set it upright
%language=Octave 					% some 'octave' instruction will appear in bold
}

\usepackage{url} 					% for trackbacking URL in the document; usage: \url{http://www.ntust.edu.tw}

% 插圖套件 graphicx
% 使用者工作流程是用 pdftex 還是 latex + dvipdfmx?
% 視情況而有不同的參數
% 這裡作自動判斷
% 參考自
% http://www.tex.ac.uk/cgi-bin/texfaq2html?label=ifpdf
%\newcommand\mydvipdfmxflow{dvipdfmx}
%\newcommand\mypdftexflow{pdftex}
%
%\ifx\pdfoutput\undefined
%  % not running pdftex
%  \usepackage[dvipdfm]{graphicx}
%  \newcommand\myworkflow{dvipdfmx}  % set the flag for hyperref
%\else
%  \ifx\pdfoutput\relax
%    % not running pdftex
%    \usepackage[dvipdfm]{graphicx}
%    \newcommand\myworkflow{dvipdfmx}  % set the flag
%  \else
%    % running pdftex, with...
%    \ifnum\pdfoutput>0
%      % ... PDF output
%      \usepackage[pdftex]{graphicx}
%      \newcommand\myworkflow{pdftex}  % set the flag
%    \else
%      %...DVI output
%      \usepackage[dvipdfm]{graphicx}
%      \newcommand\myworkflow{dvipdfmx}  % set the flag
%    \fi
%  \fi
%\fi

\usepackage{fancyhdr}  % 借用增強功能型 header 套件來擺放浮水印 
% (佔用了 central header)
% 不需要浮水印的使用者仍可利用此套件，產生所需的 header, footer
%
% 啟動 fancy header/footer 套件
\pagestyle{fancy}
\fancyhead{}  % reset left, central, right header to empty
\fancyfoot[C]{\thepage} %中間 footer 擺放頁碼
\renewcommand{\headrulewidth}{0pt} % header 的直線; 0pt 則無線

% 如果不需要任何浮水印，則請把下列介於 >>> 與 <<< 之間
% 的文字行關掉 (行首加上百分號)
%% 浮水印 >>> 
\input{watermark/ntust_watermark.tex}
%% <<< 浮水印



% global page layout
\newcommand{\mybaselinestretch}{1.5}  %行距 1.5 倍 + 20%, (約為 double space)
\renewcommand{\baselinestretch}{\mybaselinestretch}  % 論文行距預設值
\parskip=2ex  % 段落之間的間隔為兩個 x 的高度
\parindent = 24Pt  % 段首內縮由 CJK 控制，所以這裡就設成不內縮

%%%%%%%%%%%%%%%%%%%%%%%%%%%%%
%  end of preamble
%%%%%%%%%%%%%%%%%%%%%%%%%%%%%		% Basic environment settings, no need to alter 基本的環境設定 無需改變  

% Personal packages, you may add or delete it as necessary 自己需要的package，可以需要增刪
\usepackage{subfig}
\usepackage{tabularx}
\usepackage{amsfonts}
\usepackage{comment}
\usepackage{slashbox}
\usepackage{rotating}
\usepackage{booktabs,threeparttable}
\renewcommand*{\thefootnote}{\fnsymbol{footnote}}
\usepackage{pdfpages}
%\usepackage{bold-extra}

% The following package is for activating hyperlinks
%	across contents, sections, figures, and tables (e.g., \ref{} and \cite{}).
%	You can comment it if you do not need this hyperlink feature.
\usepackage{hyperref}

\begin{document}



	% 下列中文名詞的定義，如果以註解方式關閉取消，
% 則會以系統原先的預設值 (英文) 替代
% 名詞 \prechaptername 預設值為 Chapter
% 名詞 \postchaptername 預設值為空字串
% 名詞 \tablename 預設值為 Table
% 名詞 \figurename 預設值為 Figure

%\renewcommand\prechaptername{第} % 出現在每一章的開頭的「第 x 章」
%\renewcommand\postchaptername{章}
%\renewcommand{\tablename}{表} % 在文章中 table caption 會以「表 x」表示
%\renewcommand{\figurename}{圖} % 在文章中 figure caption 會以「圖 x」表示

% 下列中文名詞的定義，用於論文固定的各部分之命名 (出現於目錄與該頁標題)
\begin{comment}
\newcommand{\nameInnerCover}{教授推薦書}
\newcommand{\nameCommitteeForm}{論文口試委員審定書}
\newcommand{\nameCopyrightForm}{授權書}
\newcommand{\nameCabstract}{中文摘要}
\newcommand{\nameEabstract}{英文摘要}
\newcommand{\nameAckn}{誌謝}
\newcommand{\nameToc}{目錄}
\newcommand{\nameLot}{表目錄}
\newcommand{\nameTof}{圖目錄}
\newcommand{\nameToa}{演算法目錄}
\newcommand{\nameSlist}{符號說明}
\newcommand{\nameRef}{參考文獻}
\newcommand{\nameVita}{自傳}
\end{comment}
%下列為英文名詞的定義

\newcommand{\nameInnerCover}{Recommendation Letter}
\newcommand{\nameCommitteeForm}{Approval Letter}
\newcommand{\nameCopyrightForm}{Letter of Authority}
\newcommand{\nameCabstract}{Abstract in Chinese}
\newcommand{\nameEabstract}{Abstract in English}
\newcommand{\nameAckn}{Acknowledgements}
\newcommand{\nameToc}{Contents}
\newcommand{\nameLot}{List of Tables }
\newcommand{\nameTof}{List of Figures}
\newcommand{\nameToa}{List of Algorithms}
\newcommand{\nameSlist}{Symbol}
\newcommand{\nameRef}{References}
\newcommand{\nameVita}{Biography}	% Managing all Chinese vocabulary translations 在此檔案處定義文章中的中文名詞

	%----------------------------------------------------------------------------------------------------------------------------------------------------------
	%%% Following is the structure of front pages, main thesis, and back pages 以下是載入前頁、本文、後頁
	% Please do not change anything 此行請勿更動

	%----------------------------------------------------------------------------------------------------------------------------------------------------------
	% Front pages 前頁
	% Including: cover, title page, Chinese abstract, English abstract, 
	% 	acknowledgements, table of contents, list of figures, list of tables, etc.
	% You may comment out this front pages line while editing the main body of thesis,
	%	thus you can save much compilation time.
	% The structure of front pages for NTUST thesis is fixed,
	%	so do not alter the file ntust_frontpages.tex structure whatsoever,
	%	basically it consists my_names.text, my_cabstract.tex, my_eabstract.tex, my_ackn.tex, my_symbols.tex

	% 包括封面、書名頁、中文摘要、英文摘要、誌謝、目錄、表目錄、圖目錄、符號說明
	% 在撰寫各章草稿時，可以把此部份「關掉」，以節省無謂的編譯時間。
	% 實際內容由
	%    my_names.tex, my_cabstract.tex, my_eabstract.tex, my_ackn.tex, my_symbols.tex
	% 決定
	% ntust_frontpages.tex 此檔只提供整體架構的定義，不需更動
	% 在撰寫各章草稿時，可以把此部份「關掉」，以節省無謂的編譯時間。
	
	%
% this file is encoded in utf-8
% v1.7
% do not change the content of this file
% unless the thesis layout rule is changed
% 無須修改本檔內容，除非校方修改了
% 封面、書名頁、中文摘要、英文摘要、誌謝、目錄、表目錄、圖目錄、符號說明
% 等頁之格式
% this file is encoded in utf-8
%v1.7

% make the line spacing in effect
\renewcommand{\baselinestretch}{\mybaselinestretch}
\large % it needs a font size changing command to be effective

% default variables definitions
% 注意！！此處只是預設值，不需更改此處
% 請更改 my_names.tex 內容
\newcommand\cTitle{論文題目}
\newcommand\eTitle{MY THESIS TITLE}
\newcommand\myCname{王鐵雄}
\newcommand\myEname{Aron Wang}
\newcommand\myStudentID{M9315048}
\newcommand\advisorCnameA{南宮明博士}
\newcommand\advisorEnameA{Dr.~Ming Nangong}
\newcommand\advisorCnameB{李斯坦博士}
\newcommand\advisorEnameB{Dr.~Stein Lee}
\newcommand\advisorCnameC{徐　石博士}
\newcommand\advisorEnameC{Dr.~Sean~Hsu}
\newcommand\univCname{國立台灣科技大學}
\newcommand\univEname{National Taiwan University of science and technology}
\newcommand\deptCname{光電工程研究所}
\newcommand\fulldeptEname{Graduate School of Electro-Optical Engineering}
\newcommand\deptEname{Electro-optical Engineering}
\newcommand\collEname{College of Engineering}
\newcommand\degreeCname{碩士}
\newcommand\degreeEname{Master of Science}
\newcommand\cYear{九十四}
\newcommand\cMonth{六}
\newcommand\cDay{十}
%\newcounter{eYear}
\newcommand\eYear{2006}
\newcommand\eMonth{June}
\newcommand\ePlace{Chungli, Taoyuan, Taiwan}


 % user's names; to replace those default variable definitions
%
% this file is encoded in utf-8
% v1.7
% 填入你的論文題目、姓名等資料
% 如果題目內有必須以數學模式表示的符號，請用 \mbox{} 包住數學模式，如下範例
% 如果中文名字是單名，與姓氏之間建議以全形空白填入，如下範例
% 英文名字中的稱謂，如 Prof. 以及 Dr.，其句點之後請以不斷行空白~代替一般空白，如下範例
% 如果你的指導教授沒有如預設的三位這麼多，則請把相對應的多餘教授的中文、英文名
%    的定義以空的大括號表示
%    如，\renewcommand\advisorCnameB{}
%          \renewcommand\advisorEnameB{}
%          \renewcommand\advisorCnameC{}
%          \renewcommand\advisorEnameC{}

% 論文題目 (中文)
\renewcommand\cTitle{%我的碩士論文題目 
可驗證之安全系統的應用
}

% 論文題目 (英文)
\renewcommand\eTitle{%My Thesis Title  
A Verifiable Secure protocol in a Secure System 
%My Thesis Title  \mbox{$\cal{H}_\infty$} and \mbox{Al$_x$Ga$_{1-x}$As}
}

% 我的姓名 (中文)
\renewcommand\myCname{王大明}

% 我的姓名 (英文)
\renewcommand\myEname{Da-Ming Wang}

%我的學號
\renewcommand\myStudentID{M1915001}

% 指導教授A的姓名 (中文)
\renewcommand\advisorCnameA{陳明明 博士}

% 指導教授A的姓名 (英文)
\renewcommand\advisorEnameA{Dr.~Ming-Ming Cheng}

% 指導教授B的姓名 (中文)
\renewcommand\advisorCnameB{}

% 指導教授B的姓名 (英文)
\renewcommand\advisorEnameB{}

% 指導教授C的姓名 (中文)
\renewcommand\advisorCnameC{}

% 指導教授C的姓名 (英文)
\renewcommand\advisorEnameC{}

% 校名 (中文)
\renewcommand\univCname{國立台灣科技大學}

% 校名 (英文)
%\renewcommand\univEname{National Taiwan University of science and technology}

% 系所名 (中文)
\renewcommand\deptCname{資~訊~工~程~系}

% 系所全名 (英文)
%\renewcommand\fulldeptEname{Graduate School of Electro-Optical Engineering}

% 系所短名 (英文, 用於書名頁學位名領域)
%\renewcommand\deptEname{Electro-Optical Engineering}

% 學院英文名 (如無，則以空的大括號表示)
%\renewcommand\collEname{College of Electrical and Communication Engineering}

% 學位名 (中文)
\renewcommand\degreeCname{碩士}

% 學位名 (英文)
%\renewcommand\degreeEname{Master of Science}

% 口試年份 (中文、民國)
\renewcommand\cYear{九十八}

% 口試月份 (中文)
\renewcommand\cMonth{七} 

% 口試月份 (中文)
\renewcommand\cDay{七} 

% 口試年份 (阿拉伯數字、西元)
%\renewcommand\eYear{2009} 

% 口試月份 (英文)
%\renewcommand\eMonth{July}

% 學校所在地 (英文)
%\renewcommand\ePlace{Taipei, Taiwan}

%畢業級別；用於書背列印；若無此需要可忽略
\newcommand\GraduationClass{95}

%%%%%%%%%%%%%%%%%%%%%%
\newcommand\itsempty{}
%%%%%%%%%%%%%%%%%%%%%%%%%%%%%%%
%       ntust cover 封面
%%%%%%%%%%%%%%%%%%%%%%%%%%%%%%%
%
\begin{titlepage}
% no page number
% next page will be page 1

% aligned to the center of the page
\begin{center}
% font size (relative to 12 pt):
% \large (14pt) < \Large (18pt) < \LARGE (20pt) < \huge (24pt)< \Huge (24 pt)
%



\begin{figure}[htbp]
	\begin{minipage}[b]{5cm} 
		\raggedright
		\includegraphics[width=1.1in]{frontpages/ntust_logo.pdf}
		\label{fig:ntust_logo}
	\end{minipage}% 
	\begin{minipage}[b]{0.5\textwidth} 
	\centering
	\makebox[3cm][c]{\Huge{\univCname}}\\  %顯示中文校名
	\vspace{0.5cm}
	\makebox[3cm][c]{\Huge{\deptCname}}\\ %顯示中文系所名
	\vspace{0.5cm}
	\end{minipage}%
\\ 
\rule{16cm}{3pt}
\end{figure}
%\hfill


\vspace{1cm}
\makebox[6cm][s]{\textbf{\Huge{\degreeCname 論文}}}\\ %顯示論文種類 (中文)
%\makebox[8cm][s]{\textbf{\Huge{\degreeCname 論文（初稿）}}}\\ %顯示論文種類 (中文)
\vspace{1cm}
%
% Set the line spacing to single for the titles (to compress the lines)
\renewcommand{\baselinestretch}{1}   %行距 1 倍
%\large % it needs a font size changing command to be effective
\Large{\cTitle}\\  % 中文題目
%
\vspace{1cm}
%
\Large{\eTitle}\\ %英文題目
\vspace{5cm}
% \makebox is a text box with specified width;
% option s: stretch
% use \makebox to make sure
% 「研究生：」 與「指導教授：」occupy the same width
\hspace{4.5cm} \makebox[3cm][s]{\Large{研 究 生：}}
\Large{\myCname}  % 顯示作者中文名
\hfill \makebox[1cm][s]{}\\
%
\vspace{0.3cm}
\hspace{4.5cm} \makebox[3cm][s]{\Large{學號：}}
\Large{\myStudentID}  %顯示指導教授A中文名
\hfill \makebox[1cm][s]{}\\
%
\vspace{1cm}
\hspace{4.5cm} \makebox[3cm][s]{\Large{指導教授：}}
\Large{\advisorCnameA}  %顯示指導教授A中文名
\hfill \makebox[1cm][s]{}\\
%
% 判斷是否有共同指導的教授 B
\ifx \advisorCnameB  \itsempty
\relax % 沒有 B 教授，所以不佔版面，不印任何空白
\else
% 共同指導的教授 B
\hspace{4.5cm} \makebox[3cm][s]{}
\Large{\advisorCnameB}  %顯示指導教授B中文名
\hfill \makebox[1cm][s]{}\\
\fi
%
% 判斷是否有共同指導的教授 C
\ifx \advisorCnameC  \itsempty
\relax % 沒有 C 教授，所以不佔版面，不印任何空白
\else
% 共同指導的教授 C
\hspace{4.5cm} \makebox[3cm][s]{}
\Large{\advisorCnameC}  %顯示指導教授C中文名
\hfill \makebox[1cm][s]{}\\
\fi
%
\vfill
\makebox[10cm][s]{\Large{中華民國\cYear 年\cMonth 月\cDay 日}}%
%
\end{center}
% Resume the line spacing to the desired setting
\renewcommand{\baselinestretch}{\mybaselinestretch}   %恢復原設定
% it needs a font size changing command to be effective
% restore the font size to normal
\normalsize
\end{titlepage}
%%%%%%%%%%%%%%

%% 從摘要到本文之前的部份以小寫羅馬數字印頁碼
% 但是從「書名頁」(但不印頁碼) 就開始計算
%\setcounter{page}{1}
\pagenumbering{Roman}
%\pagenumbering{arabic}
%%%%%%%%%%%%%%%%%%%%%%%%%%%%%%%
%       指導教授推薦書 
%%%%%%%%%%%%%%%%%%%%%%%%%%%%%%%
%
% insert the printed standard form when the thesis is ready to bind
% 在口試完成後，再將已簽名的推薦書放入以便裝訂
% create an entry in table of contents for 推薦書
% 目前送出空白頁
%\newpage{\thispagestyle{empty}\addcontentsline{toc}{chapter}{\nameInnerCover}\mbox{}\clearpage}%
%\newpage

% 判斷是否要浮水印？
\ifx\mywatermark\undefined 
  \thispagestyle{empty}  % 無頁碼、無 header (無浮水印)
\else
  \thispagestyle{EmptyWaterMarkPage} % 無頁碼、有浮水印
\fi

%%%%%%%%%%%%%%%%%%%%%%%%%%%%%%%%%%%%%%%%%%%%%%%%%%%%%%%%%%%%%%%
%%no page number
%% create an entry in table of contents for 書名頁
%\addcontentsline{toc}{chapter}{\nameInnerCover}
%
%
%% aligned to the center of the page
%\begin{center}
%% font size (relative to 12 pt):
%% \large (14pt) < \Large (18pt) < \LARGE (20pt) < \huge (24pt)< \Huge (24 pt)
%% Set the line spacing to single for the titles (to compress the lines)
%\renewcommand{\baselinestretch}{1}   %行距 1 倍
%% it needs a font size changing command to be effective
%%中文題目
%\Large{\cTitle}\\ %%%%%
%\vspace{1cm}
%% 英文題目
%\Large{\eTitle}\\ %%%%%
%%\vspace{1cm}
%\vfill
%% \makebox is a text box with specified width;
%% option s: stretch
%% use \makebox to make sure
%% 「研究生：」 與「指導教授：」occupy the same width
%\makebox[3cm][s]{\large{研 究 生：}}
%\makebox[3cm][l]{\large{\myCname}} %%%%%
%\hfill
%\makebox[2cm][s]{\large{Student: }}
%\makebox[5cm][l]{\large{\myEname}}\\ %%%%%
%%
%%\vspace{1cm}
%%
%\makebox[3cm][s]{\large{指導教授：}}
%\makebox[3cm][l]{\large{\advisorCnameA}} %%%%%
%\hfill
%\makebox[2cm][s]{\large{Advisor: }}
%\makebox[5cm][l]{\large{\advisorEnameA}}\\ %%%%%
%%
%% 判斷是否有共同指導的教授 B
%\ifx \advisorCnameB  \itsempty
%\relax % 沒有 B 教授，所以不佔版面，不印任何空白
%\else
%%共同指導的教授B
%\makebox[3cm][s]{}
%\makebox[3cm][l]{\large{\advisorCnameB}} %%%%%
%\hfill
%\makebox[2cm][s]{}
%\makebox[5cm][l]{\large{\advisorEnameB}}\\ %%%%%
%\fi
%%
%% 判斷是否有共同指導的教授 C
%\ifx \advisorCnameC  \itsempty
%\relax % 沒有 C 教授，所以不佔版面，不印任何空白
%\else
%%共同指導的教授C
%\makebox[3cm][s]{}
%\makebox[3cm][l]{\large{\advisorCnameC}} %%%%%
%\hfill
%\makebox[2cm][s]{}
%\makebox[5cm][l]{\large{\advisorEnameC}}\\ %%%%%
%\fi
%%
%% Resume the line spacing to the desired setting
%\renewcommand{\baselinestretch}{\mybaselinestretch}   %恢復原設定
%\large %it needs a font size changing command to be effective
%%
%\vfill
%\makebox[4cm][s]{\large{\univCname}}\\% 校名
%\makebox[6cm][s]{\large{\deptCname}}\\% 系所名
%\makebox[3cm][s]{\large{\degreeCname 論文}}\\% 學位名
%%
%%\vspace{1cm}
%\vfill
%\large{A Thesis}\\%
%\large{Submitted to }%
%%
%\large{\fulldeptEname}\\%系所全名 (英文)
%%
%%
%\ifx \collEname  \itsempty
%\relax % 沒有學院名 (英文)，所以不佔版面，不印任何空白
%\else
%% 有學院名 (英文)
%\large{\collEname}\\% 學院名 (英文)
%\fi
%%
%\large{\univEname}\\%校名 (英文)
%%
%\large{in Partial Fulfillment of the Requirements}\\
%%
%\large{for the Degree of}\\
%%
%\large{\degreeEname}\\%學位名(英文)
%
%\large{in}\\
%%
%\large{\deptEname}\\%系所短名(英文；表明學位領域)
%%
%\large{\eMonth\ \eYear}\\%月、年 (英文)
%%
%\large{\ePlace}% 學校所在地 (英文)
%\vfill
%\large{中華民國}%
%\large{\cYear}% %%%%%
%\large{年}%
%\large{\cMonth}% %%%%%
%\large{月}\\
%\end{center}
%% restore the font size to normal
%\normalsize
%\clearpage


%%%%%%%%%%%%%%%%%%%%%%%%%%%%%%%%%%%%%%%%%%%%%%%%%%%%%%%%%%%%%%%%%%%%%
%%%%%%%%%%%%%%%%%%%%%%%%%%%%%%%
%       論文口試委員審定書 (計頁碼，但不印頁碼) 
%%%%%%%%%%%%%%%%%%%%%%%%%%%%%%%
%
% insert the printed standard form when the thesis is ready to bind
% 在口試完成後，再將已簽名的審定書放入以便裝訂
% create an entry in table of contents for 審定書
% 目前送出空白頁

%\newpage{\thispagestyle{empty}\addcontentsline{toc}{chapter}{\nameCommitteeForm}\mbox{}\clearpage}%


%%%%%%%%%%%%%%%%%%%%%%%%%%%%%%%
%       中文摘要 
%%%%%%%%%%%%%%%%%%%%%%%%%%%%%%%
%
\newpage
\thispagestyle{plain}  % 無 header，但在浮水印模式下會有浮水印
% create an entry in table of contents for 中文摘要
\addcontentsline{toc}{chapter}{\nameCabstract}

% aligned to the center of the page
\begin{center}
% font size (relative to 12 pt):
% \large (14pt) < \Large (18pt) < \LARGE (20pt) < \huge (24pt)< \Huge (24 pt)
% Set the line spacing to single for the names (to compress the lines)
\renewcommand{\baselinestretch}{1}   %行距 1 倍
% it needs a font size changing command to be effective
%\large{\cTitle}\\  %中文題目
%\vspace{0.83cm}
% \makebox is a text box with specified width;
% option s: stretch
% use \makebox to make sure
% each text field occupies the same width
%\makebox[1.5cm][s]{\large{學生：}}
%\makebox[3cm][l]{\large{\myCname}} %學生中文姓名
%\hfill
%
%\makebox[3cm][s]{\large{指導教授：}}
%\makebox[3cm][l]{\large{\advisorCnameA}} \\ %教授A中文姓名
%
% 判斷是否有共同指導的教授 B
\ifx \advisorCnameB  \itsempty
\relax % 沒有 B 教授，所以不佔版面，不印任何空白
\else
%共同指導的教授B
\makebox[1.5cm][s]{}
\makebox[3cm][l]{} %%%%%
\hfill
\makebox[3cm][s]{}
\makebox[3cm][l]{\large{\advisorCnameB}}\\ %教授B中文姓名
\fi
%
% 判斷是否有共同指導的教授 C
\ifx \advisorCnameC  \itsempty
\relax % 沒有 C 教授，所以不佔版面，不印任何空白
\else
%共同指導的教授C
\makebox[1.5cm][s]{}
\makebox[3cm][l]{} %%%%%
\hfill
\makebox[3cm][s]{}
\makebox[3cm][l]{\large{\advisorCnameC}}\\ %教授C中文姓名
\fi
%
%\vspace{0.42cm} %2013/06/06
\vspace{0.38cm}
%
%\large{\univCname}\large{\deptCname}\\ %校名系所名
%\vspace{0.83cm}
%\vfill
\makebox[2.5cm][s]{\LARGE{中文摘要}}\\
\end{center}
% Resume the line spacing to the desired setting
\renewcommand{\baselinestretch}{\mybaselinestretch}   %恢復原設定
%it needs a font size changing command to be effective
% restore the font size to normal
\normalsize
%%%%%%%%%%%%%
	
   分散式詢問及監督系統主要被用於分散式資料如檔案或是紀錄的維護上。網路內的使用者可以自系統中查詢所需資料如總和或平均值，若同時將所有原始資料做傳遞及運算，將耗費相當大的網路頻寬及運算資源。於是，內網路聚集技術被提出來降低分散式詢問及監督系統的負擔。然而，這個技術卻容易遭受安全威脅。過去的研究大多假設資料來源為可信任，並針對聚集架構進行安全性的研究。然而，我們認為聚集查詢結果應該在面對惡意攻擊者將錯誤資料置入資料串流進行聚集前，就應該具備強健的容錯能力。傳統上，一個強健的估計值被定義為即使資料來源有誤時，亦能維持一定程度正確性的聚集結果。許多常見的強健估計值是建立在有序統計學上，因此，我們將重心放在內網路計算上之有序統計的可驗證技術。此技術的挑戰為在網路遭受惡意團體介入聚集程序時，仍能確保聚集結果或近似結果的準確性。

%%%%%%%%%%%%%%%%%%%%%%%%%%%%%%%
%       英文摘要 
%%%%%%%%%%%%%%%%%%%%%%%%%%%%%%%
%
\newpage
\thispagestyle{plain}  % 無 header，但在浮水印模式下會有浮水印

% create an entry in table of contents for 英文摘要
\addcontentsline{toc}{chapter}{\nameEabstract}

% aligned to the center of the page
\begin{center}
% font size:
% \large (14pt) < \Large (18pt) < \LARGE (20pt) < \huge (24pt)< \Huge (24 pt)
% Set the line spacing to single for the names (to compress the lines)
\renewcommand{\baselinestretch}{1}   %行距 1 倍
%\large % it needs a font size changing command to be effective
%\large{\eTitle}\\  %英文題目
%\vspace{0.83cm}
% \makebox is a text box with specified width;
% option s: stretch
% use \makebox to make sure
% each text field occupies the same width
%\makebox[2cm][s]{\large{Student: }}
%\makebox[5cm][l]{\large{\myEname}} %學生英文姓名
%\hfill
%
%\makebox[2cm][s]{\large{Advisor: }}
%\makebox[5cm][l]{\large{\advisorEnameA}} \\ %教授A英文姓名
%
% 判斷是否有共同指導的教授 B
\ifx \advisorCnameB  \itsempty
\relax % 沒有 B 教授，所以不佔版面，不印任何空白
\else
%共同指導的教授B
\makebox[2cm][s]{}
\makebox[5cm][l]{} %%%%%
\hfill
\makebox[2cm][s]{}
\makebox[5cm][l]{\large{\advisorEnameB}}\\ %教授B英文姓名
\fi
%
% 判斷是否有共同指導的教授 C
\ifx \advisorCnameC  \itsempty
\relax % 沒有 C 教授，所以不佔版面，不印任何空白
\else
%共同指導的教授C
\makebox[2cm][s]{}
\makebox[5cm][l]{} %%%%%
\hfill
\makebox[2cm][s]{}
\makebox[5cm][l]{\large{\advisorEnameC}}\\ %教授C英文姓名
\fi
%
%\vspace{0.42cm}
%\large{Submitted to }\large{\fulldeptEname}\\  %英文系所全名
%
\ifx \collEname  \itsempty
\relax % 如果沒有學院名 (英文)，則不佔版面，不印任何空白
\else
% 有學院名 (英文)
%\large{\collEname}\\% 學院名 (英文)
\fi
%
%\large{\univEname}\\  %英文校名
%\vspace{0.83cm}
%\vfill
%
\large{ABSTRACT}\\
%\vspace{0.5cm}
\end{center}
% Resume the line spacing the desired setting
\renewcommand{\baselinestretch}{\mybaselinestretch}   %恢復原設定
%\large %it needs a font size changing command to be effective
% restore the font size to normal
\normalsize
%%%%%%%%%%%%%

	Distributed querying and monitoring systems have been widely studied in recent years. These systems aim to maintain data sources, such as data set or log files, and allow users to query over those data sources. When the data sources are highly related and users only care some statistic results, like the sum or the average, it is consumed to transmit all data sources via the network. To minimize the network consumption, in-network aggregation technique is proposed. However, this technique is subject to some known attacks, such as the injection attack and the pollution attack. Prior works only considered the settings that data sources are trusted while the network is not. We study the way to relax the limitation and guarantee the aggregate queries robust to malicious or faulty data sources (also called polluted data sources). 

%%%%%%%%%%%%%%%%%%%%%%%%%%%%%%%
%       誌謝 
%%%%%%%%%%%%%%%%%%%%%%%%%%%%%%%
%
% Acknowledgment
\newpage
\chapter*{\protect\makebox[5cm][s]{\nameAckn}} %\makebox{} is fragile; need protect

\addcontentsline{toc}{chapter}{\nameAckn}
\textit{alHamdulillāh}. Years of accomplishment codified on this dissertation was not easy.
It is my honor to thank all people who have all along supported me graciously.

First and foremost, I would like to express my deepest gratitude to Doctor Fulan Teng,
whose encouragement, thoughtful guidance, brilliant ideas and wholehearted supports
enabled me to develop an understanding of this research subject.
His every advice has been very valuable to my own professional growth and will not be forgotten.
I would also like to extend my gratitude to 
all committee members in my doctoral defense for their valuable comments.

It is a pleasure to thank the members of Bla Bla Bla Laboratory.

Finally, I am heartily thankful to my family
for their unconditional love, continuous supports and cheers.

%%%%%%%%%%%%%%%%%%%%%%%%%%%%%%%
%       目錄 
%%%%%%%%%%%%%%%%%%%%%%%%%%%%%%%
%
% Table of contents

\newpage


\renewcommand{\contentsname}{\protect\makebox[5cm][s]{\nameToc}}
%\makebox{} is fragile; need protect
\addcontentsline{toc}{chapter}{\nameToc}
\tableofcontents




%%%%%%%%%%%%%%%%%%%%%%%%%%%%%%%
%       圖目錄 
%%%%%%%%%%%%%%%%%%%%%%%%%%%%%%%
%
% List of Figures
\newcommand{\loflabel}{圖}
\newpage
\renewcommand{\numberline}[1]{\loflabel~#1\hspace*{1em}}
\renewcommand{\listfigurename}{\protect\makebox[5cm][s]{\nameTof}}
%\makebox{} is fragile; need protect
\addcontentsline{toc}{chapter}{\nameTof}
\listoffigures

%%%%%%%%%%%%%%%%%%%%%%%%%%%%%%%
%       表目錄 
%%%%%%%%%%%%%%%%%%%%%%%%%%%%%%%
%
% List of Tables
\newcommand{\lotlabel}{表}
\newpage
\renewcommand{\numberline}[1]{\lotlabel~#1\hspace*{1em}}
\renewcommand{\listtablename}{\protect\makebox[5cm][s]{\nameLot}}
%\makebox{} is fragile; need protect
\addcontentsline{toc}{chapter}{\nameLot}
\listoftables

%%%%%%%%%%%%%%%%%%%%%%%%%%%%%%%
%       演算法目錄 
%%%%%%%%%%%%%%%%%%%%%%%%%%%%%%%
%
% List of Figures
%\newpage
%\renewcommand{\listalgorithmname}{\protect\makebox[5cm][s]{\nameToa}}
%\makebox{} is fragile; need protect
%\addcontentsline{toc}{chapter}{\nameToa}
%\listofalgorithms


%%%%%%%%%%%%%%%%%%%%%%%%%%%%%%%
%       符號說明 
%%%%%%%%%%%%%%%%%%%%%%%%%%%%%%%
%
% Symbol list
% define new environment, based on standard description environment
% adapted from p.60~64, <<The LaTeX Companion>>, 1994, ISBN 0-201-54199-8
%\newcommand{\SymEntryLabel}[1]%
% {\makebox[3cm][l]{#1}}
%
%\newenvironment{SymEntry}
%   {\begin{list}{}%
%       {\renewcommand{\makelabel}{\SymEntryLabel}%
%        \setlength{\labelwidth}{3cm}%
%        \setlength{\leftmargin}{\labelwidth}%
%        }%
%   }%
%   {\end{list}}
%%
%\newpage
%\chapter*{\protect\makebox[5cm][s]{\nameSlist}} %\makebox{} is fragile; need protect
%\addcontentsline{toc}{chapter}{\nameSlist}
%\input{frontpages/my_symbols.tex}


\newpage
\setcounter{page}{1}
\pagenumbering{arabic}
%% 論文本體頁碼回復為阿拉伯數字計頁，並從頭起算
%\pagenumbering{arabic}
%%%%%%%%%%%%%%%%%%%%%%%%%%%%%%%% 

	%----------------------------------------------------------------------------------------------------------------------------------------------------------
	% Main body 論文主體
	% You are suggested to separate chapters into different files. 
	% 	For example: introduction.tex, related_work.tex, conclusion.tex
	%	Then put the files under sub folder "section/"

	% 建議以「章」為檔案分割的依據。
	% 以下為建議的命名分類
	%   introduction.tex, related_work.tex, conclusion.tex
	% 做為這幾個「章」的檔案名稱，並將檔案存放於資料夾 sections/ 下
	% 實際命名方式可以隨你意
	% 在撰寫各章草稿時，可以把其他章節關掉 (行首加百分號)

	%\renewcommand\thefigure{\arabic{chapter}-\arabic{figure}}
%\bibliographystyle{unsrt} 
\chapter{緒論}
\label{cha:intro} 

網路發展興盛至今，小至個人，大至政府單位與各機關組織，都相當仰賴網路的使用，但許多人仍然對資安危機意識較低，針對資訊安全產品的投資也相對較少，加上對於資訊安全軟體工具缺乏有系統的整理，以致於未能有效運用。為此，本手冊蒐集整理相關開放源碼（Open Source）的資訊安全軟體工具，並透過專業人員實際操作演練，加以彙整並集結成冊，希冀透過本手冊的幫助，不僅能給予初學者對於資安工具軟體初步認識，也讓資訊從業人員在資訊安全工具上能有更多的選擇與應用。

資安開放源碼軟體的發展，往往會公開其發展技術及運用的原理，配合程式碼的開放，使得開放源碼軟體具有相當大的彈性，並根據個人使用情況所需，進行軟體的編修與整合，以求適應各種作業環境所需。使用開放源碼軟體所需負擔的金錢成本，遠低於商業付費軟體，可降低企業組織對資訊科技產品的部分支出，不需要過度仰賴軟體製造商的技術支援與更新，也能減少相對應的軟體開發時程~\cite{GarofalakisHM07}。由於目前多數的資安開放源碼軟體的開發多為國外組織，因此較缺乏中文化介面，且部分軟體工具的使用，需要具備相當程度的專業知識~\cite{MaddenFHH02}，並非人人皆可輕易上手。本手冊擬透過中文化的工具介紹，減緩國內使用者入門的負擔。

由於現今網路環境日益複雜，遭受網路攻擊的事件層出不窮，網路安全越來越受到各界重視。網路掃描是網路安全的根本，也是攻擊者對目標主機進行攻擊的首要步驟，因此，了解網路掃描的攻擊與防禦，將有助於網路管理者提升網域的安全管理。此外，網路流量代表所有網路訊息的傳送，能提供管理者即時了解網路狀況，藉此檢視網路情況正常與否。本手冊將針對以上兩類的開放源碼軟體，逐一介紹其功能、安裝、操作與軟體評比，令讀者對相關的資訊軟體能有所了解，並進一步應用於資訊安全的監測與控管。以下即對網路掃描及流量監控兩大類軟體，進行整理與原理說明~\cite{KotiVDSD07, SubrPPKG06, ShengLMJ07, Wagner04}。
	
Insecure.org網站曾於2003年及2006年間調查各使用者喜愛的工具軟體，其中2006年收到3,243位受訪者的回覆，受訪對象涵括各界對資安工具有持續研究與發展的學者及廠商，包括Insecure.org自身、研究網際網路議題的機構、發展開放源碼軟體的組織，與其他著名的資安網站（如Open Source Security、Honeypots和IDS Focus等），並根據調查結果選出前100大網路安全工具（Top 100 Network Security Tools）。本手冊從中篩選了數套較廣泛應用且屬於開放源碼的工具軟體進行蒐集整理，以提供使用者參考學習使用。以下針對各工具軟體予以介紹相關的資訊安全基礎知識（包括專業術語解釋、專有名詞解析）、軟體安裝與使用方式，以及防駭相關知識。使用者將能透過本手冊掌握並熟悉更多相關熱門工具，且對資訊安全攻防技術有更進一步的認識。如表~\ref{tab:system}所示。 


\begin{table}[t!]
  \begin{center}
    \caption{The relation of aggregation overhead between different techniques}
    \label{tab:system}
    \begin{tabular}{|c|c c c|}
      \hline
       & Space usage & Communication & Query \\
       & of root aggregator & overhead & requirement \\
      \hline
      Traditional warehouse & $n$ & $O(n)$ & $O(n)$ \\
      \hline
      AM-FM sketch technique & $\log a$ & $O(\log n)$ &  $O(a\log n)$ \\
      \hline
      ``prototypical PHI query'' & $\log a$ & $O(\log n)$ & $O(\log n)$ \\
      \hline
      \end{tabular}
  \end{center}
\end{table}

	\renewcommand\thetable{\arabic{chapter}-\arabic{table}}
%\renewcommand\thefigure{\arabic{chapter}-\arabic{figure}}
\renewcommand{\theequation}{\arabic{chapter}-\arabic{equation}}
\chapter{研究方法}

	
	\chapter{Conclusions}
\label{cha:conclusions}


With the rapid growth in integrated circuit, digital signal processing, and other emerging technologies, people nowadays can easily purchase electronic devices, such as personal computers, laptops, cellular phones, and tablets.  
By utilizing these devices, people can communicate with each other through wireless communication and increase work performance.      
However, any malicious user may misuse these devices and launch serious attack to make illegal profit, such as identity stealing or password cracking on a bank account.
Therefore, it is essential to develop robust methods to solve the identity problems.

\section{Future Work} 
%\paragraph{Future work} 
With the rapid growth in integrated circuit, digital signal processing, and other emerging technologies, people nowadays can easily purchase electronic devices, such as personal computers, laptops, cellular phones, and tablets.  
By utilizing these devices, people can communicate with each other through wireless communication and increase work performance.      
However, any malicious user may misuse these devices and launch serious attack to make illegal profit, such as identity stealing or password cracking on a bank account.
Therefore, it is essential to develop robust methods to solve the identity problems.



	%----------------------------------------------------------------------------------------------------------------------------------------------------------
	% Back pages 後頁
	% Including: references, appendixes, and curriculum vitae
	% The structure of back pages for NTUST thesis is fixed,
	%	so do not alter the file ntust_backpages.tex structure whatsoever.
	% You can also comment out this back pages line while editing the main body of thesis.

	% 包括參考文獻、附錄、自傳
	% 實際內容由
	%    my_bib.bib, my_appendix.tex, my_vita.tex
	% 決定
	% ntust_backpages.tex 此檔只提供整體架構的定義，不需更動
	% 在撰寫各章草稿時，可以把此部份「關掉」，以節省無謂的編譯時間。

	%
% This file is encoded in utf-8
% v1.7
%

%%% References Document/ Bibliography 
\newpage
\addcontentsline{toc}{chapter}{\nameRef}
\renewcommand{\bibname}{\protect\makebox[5cm][s]{\nameRef}}		% '\makebox{}' is fragile, '\protect' will help
\bibliographystyle{ieeetr}										% Using IEEE Trans journal format
{\footnotesize
\bibliography{my_bib}  											% 'my_bib' is a BibTeX file located in the parent directory, it comprises all the references for the thesis
}

%%% Appendix (uncomment the following line to include the appendixes)
%%
% This file is encoded in utf-8
% v1.7
%

%%% Please copy the following TeX as the header code of each appendix (appendix 1, appendix 2, etc.):

%%% Header code begins >>>

\newpage
\chapter*{Appendix 1：MATLAB Source Code List} 							% Please change the values inside the curly brackets (the appendix number and the title)
\addcontentsline{toc}{chapter}{Appendix 1：MATLAB Source Code List} 	% It is recommended to set the same title as above
\renewcommand{\thechapter}{1} 											% if it is appendix 2, then it should be written as {2}

\setcounter{equation}{0} 
\setcounter{figure}{0} 
\setcounter{footnote}{0} 
\setcounter{section}{0} 
\setcounter{subsection}{0}
\setcounter{subsubsection}{0}
\setcounter{table}{0} 

%%% <<< Header code ends

% Appendix content is started from here

\lstinputlisting{example/example_prog_list.m}

%%% For generating appendix 2, 3, and so on, please continue by incrementing the order number in the TeX above
% Every appendix will automatically start as a new page

%%% Autobiography/ Vitae
%\newpage
%\chapter*{\protect\makebox[5cm][s]{\nameVitae}} 				% '\makebox{}' is fragile, '\protect' will help
%\addcontentsline{toc}{chapter}{\nameVitae}
%Fulan Chen is currently a Ph.D candidate in Department of Computer Science and Information Engineering from National Taiwan University of Science and Technology, Taipei, Taiwan. He received M.S degree in Department of Computer Science and Information Engineering from National Taiwan University of Science and Technology in 2015. His current research interests include the Human-Computer Interaction, Tele-operation and Robotics. 



%%%%%%%%%%%%%%%%%%%%%%%%%%%%%%%%%%%%%%%%%%%%%%%%%%%%%%%%%%%%%%
%   Letter of Authorization (LoA/ 授權書)
%     ... the page numbers will be counted but not be printed
%%%%%%%%%%%%%%%%%%%%%%%%%%%%%%%%%%%%%%%%%%%%%%%%%%%%%%%%%%%%%%
%
% Insert the printed standard form when the thesis is ready to bind
% After the oral defense, put the LoA (signed by advisor) into the thesis, then bind it
% Create an entry in table of contents for this LoA
% Currently only generate an empty page

\newpage{\thispagestyle{empty}\addcontentsline{toc}{chapter}{\nameCopyrightForm}\mbox{}\clearpage}

%\includepdf{backpages/a.pdf}{\addcontentsline{toc}{chapter}{\nameCopyrightForm}}


%%%%% useless %%%%%
%\begin{figure}[!t]
%\centering
%\includegraphics[width=\paperwidth]{backpages/LetterofAuthority.pdf}
%\end{figure}
%\includepdf[addtotoc={1,chapter,-1,\nameCopyrightForm,letter}]{backpages/LetterofAuthority.pdf}
%\includepdf{backpages/LetterofAuthority.pdf}
%\addcontentsline{toc}{chapter}{\nameCopyrightForm}\mbox{}

\end{document} 
 
